\chapter*{Abstract}

Goal-oriented remote meetings often suffer from unequal participation, incomplete information exchange, and
premature consensus, which can impair decision-making. While existing AI-based meeting tools primarily
support preparation and documentation, the potential of proactive AI moderation during ongoing discussions
remains insufficiently explored. This thesis investigates how a proactive, context-sensitive AI system can
support convergent decision-making in remote meetings.

First, a literature review and three expert interviews were conducted to identify relevant facilitation
strategies and challenges. Based on these findings, a prototype was developed to address hidden-profile
scenarios, in which relevant information is distributed asymmetrically among participants. The system aims to
encourage balanced participation, elicit unshared information, and preserve dissent without influencing the
decision itself. Its evaluation comprised three group sessions with twelve participants, using a hiring
simulation and questionnaires to assess meeting quality and perceptions of the AI moderator.

The results indicate that proactive AI moderation has potential, but its effectiveness depends on the
relevance and timing of interventions. Although participants generally perceived discussion quality
positively, the AI was often considered disruptive due to excessive, delayed, or contextually inappropriate
contributions. Nevertheless, interventions encouraging quieter participants to contribute or requesting
additional perspectives were sometimes perceived as helpful. The findings highlight that effective AI
facilitation requires not only relevant contributions but also context-sensitive and selective intervention.
Further research should focus on improving contextual awareness, intervention timing, and participant control.

Keywords: proactive AI moderation, remote meetings, group decision-making, meeting facilitation, human-AI interaction